Large language models are increasingly used to triage and resolve consumer complaints, and the way complainants write varies systematically with age, education and socioeconomic status (SES). We audit three remedy-decision systems using sociolinguistic rewrites of US consumer financial complaints that are instructed to preserve every fact: two generative LLMs that return a single label (GPT-4o-mini and GPT-5-mini), and a probabilistic ``System One'' model that returns a distribution over remedies (Jev). Using pre-registered hypotheses tested once on held-out complaints, we separate two phenomena that single-label audits conflate. At the group scale, older, college-educated and higher-SES writing styles receive slightly more generous remedies. These profiling effects are reliable and replicate, and they are small in absolute terms: on held-out complaints every monetary-remedy gap is within 0.5 percentage points. At the individual scale, two style versions of the same complaint receive different monetary decisions 2.6--6.3\% of the time. In a variance-component analysis of the probabilistic model, persona components shared across the two rewrite models account for an estimated 1.6\% of within-complaint variation; most of the rest comes from which model wrote the rewrite, guise-specific paraphrase, and option order and inference noise. The persona gaps are statistically associated mainly with delivery (vocabulary, grammar and politeness) rather than content, and are larger where the style cue is recognisable. Positive controls show that all three systems respond 49--387× more strongly to remedy-relevant facts than to style. Graded probabilistic outputs give smoother scores but, at matched monetary volume, barely change how often monetary decisions disagree. In a pre-registered mitigation, rewriting every complaint into one standard style reduced GPT-4o-mini's education and SES gaps by 94\% and 87\% in point estimates while retaining 96--102\% of fact sensitivity; it reversed two of the probabilistic model's gaps and did not make individual decisions more stable. In an independent audit, two LLMs from outside the audited systems flagged material fact changes in 3.3\% and 11.7\% of rewrites, and instability persisted among rewrites that neither flagged. We argue that fairness audits of LLM decision systems should report group disparities and individual instability as separate outcomes, evaluate both on the decisions people receive, and audit any mitigation with the same design.
